\input{estilo}
% Propuesta: el sharding aleatorio en tres paneles iguales —el frontend y los
% tres shards en el mismo lugar— con las tres operaciones. Send: un shard al
% azar. Receive: también al azar; si cae en otro shard vuelve vacío (dos de
% cada tres veces), y cuando encuentra m₁ lo devuelve con un handle opaco que
% dice dónde estaba, escrito en palabras arriba del panel. Delete: el handle lleva al shard correcto. Agregado del
% texto: los tres paneles, el handle y los pies.
\begin{document}
\begin{figura}\begin{tikzpicture}[
    sh/.style={caja, minimum width=10mm, minimum height=9mm},
    fe/.style={caja, minimum width=14mm, minimum height=9mm},
  ]
  \newcommand\msj[1]{\node[draw=segundo, fill=segundo!14, text=segundo, minimum size=4.6mm,
    inner sep=0pt, font=\scriptsize] at #1 {\var{m}\textsubscript{1}};}
  % un panel: shards con su base, frontend arriba; \panel{x}{título}{pie}
  \newcommand\panel[3]{\begin{scope}[xshift=#1cm]
    \node[fe] (f) at (2.0,2.0) {FE};
    \foreach \i/\x in {1/0, 2/2.0, 3/4.0} {
      \node[sh] (b\i) at (\x,0) {\var{B}\textsubscript{\i}};
      \pic at (\x,-2.05) {base};
      \draw (b\i.south) -- (\x,-0.95);
    }
    \node at (2.0,-2.7) {#2};
    \node[etiq, text=suave] at (2.0,-3.2) {#3};
  \end{scope}}

  % send
  \panel{0}{send}{un shard al azar}
  \draw[msg] (2.0,1.55) -- (0,0.45);
  \msj{(0,-1.45)}

  % receive
  \panel{6.4}{receive}{al azar: dos de cada tres, vacío}
  \draw[msg, draw=suave] (8.4,1.55) -- (10.4,0.45);
  \node[etiq, text=acento, anchor=west] at (9.8,1.15) {vacío};
  \draw[msg] (8.4,1.55) -- (6.4,0.45);
  \msj{(6.4,-1.45)}
  \draw[respuesta] (8.4,2.45) -- ++(0,0.5);
  \node[anchor=south, align=center] at (8.4,2.95) {\var{m}\textsubscript{1} + handle\\[-1pt]{\scriptsize\color{suave}el handle: en qué shard quedó, opaco}};

  % delete
  \panel{12.8}{delete(handle)}{el handle dice en qué shard está}
  \draw[msg] (14.8,2.95) -- ++(0,-0.5);
  \node[anchor=south] at (14.8,2.95) {handle};
  \draw[msg] (14.8,1.55) -- (12.8,0.45);
  \node[etiq, text=segundo, font=\scriptsize] at (12.8,-1.45) {borrado};
\end{tikzpicture}\end{figura}
\end{document}
